\scanpage{439}
\mspage{6}
% archive sheet number: 438

or le \inserted{$T$-}foncteur canonique

\[
H\longrightarrow\Pi(H/T)\qquad
\substack{\text{(= cat. sur }T\text{ définie}\\
\text{par }\ t\longmapsto\Pi({}_{t}\backslash H)}
\]

est une \deleted{\unsure{$S$}} équivalence, puisque $H$ est un \deleted{groupoïde} fibré en groupoïdes sur $T$. Donc on peut \deleted{\unsure{l'inverser en}} trouver un \inserted{$T$-}quasi-inverse

\[
\Pi(H/T)\longrightarrow H
\]

et en composant on trouve

\[
f_!(f^*H)\longrightarrow\Pi(H/T)\unsure{\mathrel{-}|\longrightarrow}H
\]

la flèche cherchée.

Ayant les \inserted{\unsure{2|2\textsuperscript{e}}} flèches d'adjonction, je m'arrête là dans les vérifications !

\unsure{Questions|Question} de changement de base

\diagram{\begin{tikzcd}[ampersand replacement=\&,column sep=8mm,row sep=6mm]
\& G \& {} \& {}\\
\& S\arrow[d,"f"'] \& S'\arrow[l,"q"']\arrow[d,"f'"] \& {}\\
(*) \& T \& T'\arrow[l,"p"] \& p^*f_!(G)\longleftarrow f'_!q^*(G)
\arrow[from=3-3,to=2-2,shorten <=5mm,shorten >=6mm,"\unsure{\gamma|r}"]
\end{tikzcd}}

La flèche est définie, pourvu qu'on ait commutativité \inserted{\unsure{ou} une flèche de commutativité $\gamma\colon pf'\unsure{\mathrel{-}|\longrightarrow}fq$ (pas néc. iso).} \deleted{\unsure{Par}} Par adjonction, il revient au \unsure{$=$|même} de donner

\diagram{\begin{tikzcd}[ampersand replacement=\&,column sep=11mm,row sep=5mm]
p_!f'_!q^*G\arrow[r,dash]\arrow[d,"\text{transitivité}"',"\simeq",dash] \& f_!(G)\\
(pf')_!q^*G\arrow[d,"\gamma"] \& {}\\
(fq)_!q^*G\arrow[uur,"\text{car }q_!q^*G\longrightarrow G"'] \& {}\\
f_!(q_!q^*)G\arrow[u,"\text{transitivité}"',"\simeq",dash] \& {}
\end{tikzcd}}

On se demande quand la flèche est iso.
